\chapter{Contrats à terme sur obligations}\label{ch:m2:bond-futures}

Le contrat à terme sur notes du Trésor à dix ans du Chicago Board of Trade promet la livraison
d'une note qui n'existe pas. Ses modalités
décrivent un panier~: toute note du Trésor dont la durée restante, au début du mois de livraison, est comprise
entre six ans et demi et huit ans peut être livrée, et le vendeur, dit le short, choisit
laquelle. Pour rendre interchangeables des notes de coupons et d'échéances différents, chacune
est multipliée par un facteur publié par la bourse~; ce facteur n'est exact que lorsque les rendements
valent 6~\%, ce qui est rare, si bien qu'une note du panier est toujours moins chère à livrer
que les autres, et le contrat à terme se comporte comme un terme sur cette note,
augmenté de la valeur du droit du vendeur de changer d'avis. Les
hedge funds détiennent des centaines de milliards de
dollars de notes adossées à des ventes à terme pour gagner la faible différence entre
les deux. Ce chapitre évalue le contrat, identifie la note sur laquelle il est réellement
adossé et valorise ce choix.

\section{Le contrat et son panier}

\begin{definition}[Panier livrable]\label{def:m2:bond-futures:basket}
Le \emph{panier livrable}\index{panier livrable} d'un contrat à terme sur obligations est
l'ensemble des titres que le short peut livrer, défini par les règles de la bourse
selon l'échéance restante (et parfois initiale) au début du
mois de livraison.
\end{definition}

\begin{dated}{2026-09}{Deux contrats à dix ans}\label{dat:m2:bond-futures:contracts}
\textbf{Contrat CBOT à dix ans} (le contrat «~6.5 to 8-year~»,
code \texttt{ZN})~: USD~100\,000 de valeur faciale~; tick égal à la moitié d'un 32e, USD~15.625~; notes livrables
d'une durée initiale d'au plus dix ans et d'une durée restante d'au moins
six ans et six mois et moins de huit ans~; à partir du contrat de mars 2026,
les notes rouvertes sont également livrables. Dernier jour de négociation~: le
septième jour ouvré précédant le dernier jour ouvré du mois de livraison~;
livraison jusqu'au dernier jour ouvré. \textbf{Contrat à terme Eurex Euro-Bund}
(\texttt{FGBL})~: EUR~100\,000 d'une obligation notionnelle à 6~\%~; obligations
fédérales allemandes livrables, d'une durée restante de 8.5 à 10.5 ans.
\end{dated}

\section{Facteurs de concordance et prix de facturation}

\begin{definition}[Facteur de concordance et prix de facturation]\label{def:m2:bond-futures:cf}
Le \emph{facteur de concordance}\index{facteur de concordance} d'une note livrable
est, approximativement, le prix d'une unité de pair à un rendement de 6~\%, calculé par
la formule de la bourse, avec le coupon de la note arrondi au huitième le plus proche
et sa durée restante, comptée à partir du premier jour du mois de livraison et arrondie par défaut,
à un trimestre entier pour les contrats à dix ans et sur obligations. Un short qui livre
la note reçoit le \emph{prix de facturation}\index{prix de facturation}
\[
\text{facturation} \;=\; F \times \mathrm{CF} \;+\; \text{coupon couru à la livraison},
\]
pour 100 de valeur faciale, où $F$ est le prix de règlement-livraison du contrat à terme.
\end{definition}

\begin{method}[Le facteur de concordance de la bourse]\label{met:m2:bond-futures:cf}
Avec le coupon $c$ exprimé en décimales, $n$ années entières et $z$ mois entiers depuis le
premier jour du mois de livraison jusqu'à l'échéance ($z$ arrondi par défaut à un multiple de
3 pour les contrats à dix ans et sur obligations), $v = z$ si $z < 7$ et $v = 3$
sinon (pour ces contrats), on pose $a = 1.03^{-v/6}$, $b = \frac c2
\frac{6 - v}{6}$, $c' = 1.03^{-2n}$ si $z < 7$ et $1.03^{-(2n+1)}$ sinon,
$d = \frac{c}{0.06}(1 - c')$, et
\[
\mathrm{CF} \;=\; a\,\Bigl[\frac c2 + c' + d\Bigr] - b,
\]
arrondi à quatre décimales. C'est le prix à 6~\% d'une note dont le prochain coupon
est à $v$ mois, diminué de la part de ce coupon déjà courue.
\end{method}

\begin{example}[La vérification propre à la bourse]\label{ex:m2:bond-futures:cf}
Pour la note de 3.75~\% du 15 novembre 2018 et le contrat de décembre 2008,
$n = 9$, les 11 mois restants sont arrondis par défaut à $z = 9$, donc $v =
3$~: $a = 0.985329$, $b = 0.009375$, $c' = 0.570286$, $d = 0.268571$ et
$\mathrm{CF} = 0.8357$, le facteur publié par la bourse. La construction
du chapitre reproduit les cinq exemples résolus de la bourse à la quatrième décimale.
\end{example}

Le facteur donne à une note à fort coupon une valeur proportionnellement plus élevée
à la livraison, de sorte que toutes les notes du panier seraient également intéressantes à livrer
si tous rapportaient 6~\%. À tout autre rendement, elles ne le sont pas. Lorsque les rendements sont
inférieurs à 6~\%, une note à longue duration gagne davantage que ne le permet le facteur, si bien que le short
préfère livrer la note de \emph{plus faible} duration~; lorsque les rendements sont supérieurs à 6~\%
il préfère la note de duration la plus élevée.

\section{Le moins cher à livrer}

\begin{definition}[Moins cher à livrer]\label{def:m2:bond-futures:ctd}
Le \emph{moins cher à livrer}\index{moins cher à livrer} (CTD) est la note du panier
qui coûte le moins au short à acheter et à livrer~: celle dont le prix est le plus bas
par rapport à ce qui lui sera facturé, ou, de manière équivalente, le
\omterm{def:m2:bond-futures:irr}{taux de repo implicite} le plus élevé (la \cref{def:m2:bond-futures:irr}).
\end{definition}

\begin{proposition}[À la livraison]\label{prop:m2:bond-futures:atdelivery}
Le dernier jour de livraison, si aucun arbitrage n'est possible, $F = \min_i
P_i/\mathrm{CF}_i$, où $P_i$ est le prix pied de coupon de la note $i$~: le prix du contrat à terme
converge vers le prix de la note la moins chère, divisé par son \omterm{def:m2:bond-futures:cf}{facteur de concordance}.
\end{proposition}
\begin{proof}
Si $F > P_i/\mathrm{CF}_i$ pour un certain $i$, vendre le contrat à terme, acheter la note $i$ et
la livrer~: la facturation $F\,\mathrm{CF}_i + A_i$ dépasse le coût $P_i + A_i$. Si
$F < \min_i P_i/\mathrm{CF}_i$, acheter le contrat à terme~: le short doit livrer une certaine note
$j$ valant $P_j + A_j > F\,\mathrm{CF}_j + A_j$, que le long peut revendre
plus cher que ce qu'il a payé.
\end{proof}

Avant la livraison, le contrat à terme est évalué à partir du terme sur le CTD, soit le prix auquel
acheter la note aujourd'hui et la financer en repo atteint le point mort à la livraison, et
il se négocie un peu en dessous, parce que le short détient les
options de la \cref{sec:m2:bond-futures:options}.

\section{Base brute, base nette et taux de repo implicite}

\begin{definition}[Base brute, base nette, taux de repo implicite, basis trade]\label{def:m2:bond-futures:irr}
Pour une note livrable de prix pied de coupon $P$, de facteur $\mathrm{CF}$ et de
prix à terme $F$~: la \emph{base brute}\index{base brute} (gross basis) est $P - F\,\mathrm{CF}$~;
la \emph{base nette}\index{base nette} (net basis) est la base brute diminuée du carry de la note
jusqu'à la livraison (\cref{prop:m2:repo-and-specials:carry})~; le \emph{taux de repo
implicite}\index{taux de repo implicite} (implied repo rate) est le taux de financement auquel l'achat de la note,
son financement et sa livraison à terme atteignent exactement le point mort,
\[
\text{IRR} \;=\; \frac{F\,\mathrm{CF} + A_{\text{del}} + C - (P + A_0)}{P + A_0}\,\frac{360}{d},
\]
avec le coupon couru $A_0$ aujourd'hui et $A_{\text{del}}$ à la livraison, les coupons
$C$ perçus entre-temps et $d$ jours. Un \emph{basis trade}\index{basis trade}
achète la note et vend le contrat à terme dans le rapport du facteur (position longue sur la
base), ou l'inverse (position courte sur la base).
\end{definition}

Le \omterm{def:m2:bond-futures:irr}{taux de repo implicite} est la cotation que fixe le marché à terme pour le financement~: un trader
qui détient une position longue sur la base perçoit ce taux sur le prix de la note et paie le \omterm{def:m2:repo-and-specials:rate}{taux de repo} effectif,
de sorte que l'écart constitue le gain de la transaction si la note est livrée. Pour le
CTD, il est en général un peu inférieur au taux du collatéral général (GC), et l'écart, comme la \omterm{def:m2:bond-futures:irr}{base nette},
valorise les options du short.

\begin{example}[Un panier de décembre]\label{ex:m2:bond-futures:basket}
Le 25 septembre 2026, le contrat à terme de décembre sur la note à dix ans s'échange à 111-24+
(111.765625)~; la livraison est supposée avoir lieu le 31 décembre, dans 97 jours, et le \omterm{def:m2:repo-and-specials:rate}{taux de repo} est
3.90~\%. Le tableau ci-dessous présente cinq notes livrables. La
note 3.875~\% d'août 2033, la plus courte, présente la plus petite \omterm{def:m2:bond-futures:irr}{base nette},
deux tiers d'un 32e, et le \omterm{def:m2:bond-futures:irr}{taux de repo implicite} le plus élevé, 3.82~\%, huit points de base
sous le \omterm{def:m2:repo-and-specials:rate}{taux de repo}~: c'est le CTD. La note la plus longue serait livrée à un \omterm{def:m2:bond-futures:irr}{taux de repo
implicite} de 0.05~\%. Ces notes et ces prix sont illustratifs.
\end{example}

\begin{center}
\small
\begin{tabular}{lrrrrr}
\toprule
Note & Prix pied de coupon & Facteur & \omterm{def:m2:bond-futures:irr}{Base brute} & \omterm{def:m2:bond-futures:irr}{Base nette} & \omterm{def:m2:bond-futures:irr}{Taux de repo implicite} \\
\midrule
3.875~\% août-33 & 99.132 & 0.8870 & $-0.004$ & 0.020 & 3.82~\% \\
4.125~\% nov.-33 & 100.549 & 0.8971 & 0.284 & 0.261 & 2.95~\% \\
4.25~\% févr.-34 & 101.262 & 0.9012 & 0.539 & 0.488 & 2.12~\% \\
4.625~\% mai-34 & 103.641 & 0.9200 & 0.817 & 0.695 & 1.45~\% \\
4~\% août-34 & 99.462 & 0.8806 & 1.041 & 1.036 & 0.05~\% \\
\bottomrule
\end{tabular}
\end{center}

\begin{omfigure}
\begin{tikzpicture}[font=\small,
  box/.style={draw,thick,rounded corners=2pt,align=center,minimum height=0.95cm,minimum width=2.6cm,inner sep=4pt},
  arr/.style={-{Stealth[length=2.5mm]},thick}]
\node[box,draw=omThm,fill=omThm!8] (fund) at (0,0) {basis trader\\(position longue sur la base)};
\node[box,draw=omDef,fill=omDef!8] (cash) at (-4.6,1.6) {marché au comptant};
\node[box,draw=omProp,fill=omProp!8] (repo) at (-4.6,-1.6) {prêteur en repo};
\node[box,draw=gray,fill=gray!8] (fut) at (4.6,0) {bourse de contrats à terme\\(position courte)};
\draw[arr] (cash) -- node[above=11pt,sloped,font=\footnotesize] {achète le CTD} (fund);
\draw[arr] (repo) -- node[below=11pt,sloped,font=\footnotesize] {la finance} (fund);
\draw[arr] (fund.north east) to[bend left=25] node[above=5pt,font=\footnotesize] {livre la note} (fut.north west);
\draw[arr] (fut.south west) to[bend left=25] node[below=5pt,font=\footnotesize] {prix de facturation $F\,\mathrm{CF} + A$} (fund.south east);
\end{tikzpicture}

\omcaption{Le \omterm{def:m2:bond-futures:irr}{basis trade} cash-and-carry. Le trader achète la note la moins chère,
la finance en repo, vend des contrats à terme dans le rapport du facteur et livre
la note à l'échéance contre le \omterm{def:m2:bond-futures:cf}{prix de facturation}. Il perçoit le \omterm{def:m2:bond-futures:irr}{taux de repo implicite}
et paie le \omterm{def:m2:repo-and-specials:rate}{taux de repo}~; avec l'effet de levier du repo, quelques points de base sur une
position importante, c'est là le métier.}\label{fig:m2:bond-futures:trade}
\end{omfigure}

\begin{dated}{2026-09}{L'ampleur du basis trade}\label{dat:m2:bond-futures:basis}
Une note de la Federal Reserve de juin 2026 évaluait les positions de base agrégées sur le
cash et les contrats à terme sur Treasuries à environ 830~milliards de dollars en septembre 2025,
à peu près le double du sommet du début 2020 et 3.5~\% des Treasuries détenues par le secteur privé, en valeur de marché.
Les estimations publiées vont de 350~milliards de dollars à 1.5~billion de dollars~; les
avoirs en Treasuries au comptant des hedge funds ont atteint 2~billions de dollars à la fin de 2025.
\end{dated}

\section{Les options de livraison}\label{sec:m2:bond-futures:options}

\begin{definition}[Option de qualité et option wildcard]\label{def:m2:bond-futures:options}
Le droit du short de choisir la note à livrer est l'\emph{option de
qualité}\index{option de qualité} (quality option). L'\emph{option wildcard}\index{option wildcard}
(wildcard option) est son droit, pendant la période de livraison, de décider s'il livre
après la fixation du cours de règlement du contrat à terme pour la journée, tandis que les notes
sont encore négociées~; après le dernier jour de négociation, le \omterm{def:m2:bond-futures:cf}{prix de facturation} est fixé et le
short peut livrer un quelconque des jours restants. Le droit de choisir le jour
dans le mois de livraison est l'\emph{option de calendrier} (timing option).
\end{definition}

\begin{omfigure}
\begin{tikzpicture}[font=\footnotesize]
\draw[thick,-{Stealth[length=2.5mm]}] (0,0) -- (13.4,0) node[right]{temps};
\draw[omDef,line width=3pt] (0.6,0.35) -- (9.0,0.35);
\foreach \x in {2.0,3.5,5.0,6.5,8.0}{\draw[omProp,thick] (\x,0.5) -- (\x,0.8);}
\node[anchor=south,text=omProp] at (5.0,0.85) {fenêtres wildcard quotidiennes};
\node[anchor=south,text=omDef] at (4.8,1.35) {négociation du contrat à terme~; cours de règlement fixé chaque jour};
\draw[omThm,line width=3pt] (9.0,0.35) -- (12.6,0.35);
\node[anchor=south,text=omThm,align=center] at (10.8,0.5) {prix de facturation fixé~:\\livraison possible tout jour};
\foreach \x/\t in {0.6/premier jour du mois,9.0/dernier jour de négociation,12.6/dernier jour de livraison}{
  \draw[thick] (\x,0.12) -- (\x,-0.12); \node[below=3pt] at (\x,-0.12) {\t};}
\end{tikzpicture}

\omcaption{Le mois de livraison du contrat à terme sur la note à dix ans. Le short choisit
la note (\omterm{def:m2:bond-futures:options}{option de qualité}), le jour (option de calendrier) et, chaque jour jusqu'à
la fin de la négociation, s'il livre après la publication du cours de règlement du jour
(\omterm{def:m2:bond-futures:options}{option wildcard}). La négociation cesse le septième jour ouvré précédant le dernier
jour ouvré du mois~; le \omterm{def:m2:bond-futures:cf}{prix de facturation} est alors fixé, et la livraison peut
avoir lieu tout jour ouvré restant jusqu'au dernier.}\label{fig:m2:bond-futures:timeline}
\end{omfigure}

Les trois appartiennent au short, si bien que le contrat à terme vaut moins que le terme du
CTD diminué de leur valeur, et un trader en position longue sur la base (achat de la note, vente du
contrat à terme) les détient. L'\omterm{def:m2:bond-futures:options}{option de qualité} vaut le plus lorsque le CTD pourrait
basculer d'une note à l'autre~: lorsque les prix par facteur de deux notes sont proches, ou lorsque les
taux sont assez volatils pour rendre ce basculement envisageable
(la \cref{fig:m2:bond-futures:switch}).

\begin{omfigure}
\begin{tikzpicture}
\begin{axis}[width=11cm,height=5cm,
  xlabel={variation parallèle des rendements (pb)},ylabel={prix/facteur moins le minimum},
  tick label style={font=\small},label style={font=\small},
  xmin=-200,xmax=250,ymin=-0.05,ymax=1.6,grid=major,grid style={gray!25},
  legend columns=3,legend style={font=\footnotesize,at={(0.5,-0.3)},anchor=north,draw=none,fill=none,/tikz/every even column/.append style={column sep=6pt}},
  table/col sep=comma]
\addplot[omDef,very thick] table[x=shift,y=d0]{figdata/markets-2/06-bond-futures/switch.csv};
\addplot[omProp,thick,dashed] table[x=shift,y=d1]{figdata/markets-2/06-bond-futures/switch.csv};
\addplot[gray,thick,dotted] table[x=shift,y=d2]{figdata/markets-2/06-bond-futures/switch.csv};
\addplot[omProp!60,thick] table[x=shift,y=d3]{figdata/markets-2/06-bond-futures/switch.csv};
\addplot[omThm,very thick] table[x=shift,y=d4]{figdata/markets-2/06-bond-futures/switch.csv};
\draw[omThm,thick,dashed] (axis cs:147,-0.05) -- (axis cs:147,1.6);
\node[font=\footnotesize,anchor=south east,text=omThm] at (axis cs:150,1.25) {basculement~: +147 pb};
\legend{août-33 (CTD aujourd'hui),nov.-33,févr.-34,mai-34,août-34}
\end{axis}
\end{tikzpicture}

\omcaption{À la livraison, le prix pied de coupon de chaque note divisé par son facteur, moins
le plus bas, lorsque tous les rendements évoluent ensemble. La note zéro-coupon est le CTD. À partir des
rendements d'aujourd'hui, la note la plus courte reste la moins chère jusqu'à ce que les rendements aient monté de 147
points de base, soit environ 5.5~\%~; au-delà, la note la plus longue prend le relais. Les notes et
les rendements sont donnés à titre illustratif. Données~: le tutoriel du chapitre.}\label{fig:m2:bond-futures:switch}
\end{omfigure}

\begin{omfigure}
\begin{tikzpicture}
\begin{axis}[width=11cm,height=4.6cm,ybar,bar width=20pt,
  ylabel={repo implicite (\%)},
  tick label style={font=\small},label style={font=\small},
  xmin=-0.6,xmax=4.6,ymin=0,ymax=4.6,
  xtick=data,xticklabels={août-33,nov.-33,févr.-34,mai-34,août-34},
  xticklabel style={font=\footnotesize},
  table/col sep=comma]
\addplot[fill=omDef!55,draw=omDef] table[x=i,y=irr]{figdata/markets-2/06-bond-futures/irr.csv};
\draw[omThm,thick,dashed] (axis cs:-0.6,3.90) -- (axis cs:4.6,3.90);
\node[font=\footnotesize,anchor=south east,text=omThm] at (axis cs:4.55,3.92) {taux de repo 3.90~\%};
\end{axis}
\end{tikzpicture}

\omcaption{\omterm{def:m2:bond-futures:irr}{Taux de repo implicites} des cinq notes de
l'\cref{ex:m2:bond-futures:basket}. Celui du CTD est le plus élevé, juste sous le
\omterm{def:m2:repo-and-specials:rate}{taux de repo}~; les autres perdraient de l'argent si on les achetait puis les livrait. Données~: le
tutoriel du chapitre.}\label{fig:m2:bond-futures:irr}
\end{omfigure}

\section{Tutoriel~: trouver le moins cher à livrer}\label{tut:m2:bond-futures:ctd}

\begin{tutorial}
\textbf{Objectif.} Calculer les facteurs d'un panier, ses \omterm{def:m2:bond-futures:irr}{bases brute} et nette et
les \omterm{def:m2:bond-futures:irr}{taux de repo implicites}, trouver le CTD et localiser le rendement auquel il bascule.
\textbf{État final~:} le tableau de l'\cref{ex:m2:bond-futures:basket} et
les \cref{fig:m2:bond-futures:switch,fig:m2:bond-futures:irr}.
\begin{tutsteps}
\item \textbf{Le facteur}, à partir de la formule de la bourse.
\omcode{code/firm/ctd/firm_ctd.py}{23}{36}{Le facteur de concordance.}\label{lst:m2:bond-futures:cf}
Vérifier les exemples publiés par la bourse~: 0.9229, 0.8747, 0.8653, 0.8357, 0.7943.
\item \textbf{Base et carry.}
\omcode{code/firm/ctd/firm_ctd.py}{53}{64}{Base brute, carry jusqu'à la livraison et base nette.}\label{lst:m2:bond-futures:basis}
\item \textbf{Repo implicite et CTD.}
\omcode{code/firm/ctd/firm_ctd.py}{67}{76}{Taux de repo implicite, et la note qui le maximise.}\label{lst:m2:bond-futures:irr}
\item \textbf{Le basculement.} \texttt{futures\_demo.switch\_shift()} parcourt
des déplacements parallèles de la courbe et signale le CTD qui change à $+147$ points de base~;
\texttt{fig\_futures.py} écrit les deux graphiques.
\end{tutsteps}
\textbf{Pour aller plus loin.} Accentuer la pente de la courbe en ajoutant un point de base par
année de maturité au rendement de chaque note, et trouver le nouveau point de basculement~;
puis calculer la valeur de l'option de basculement pour plusieurs volatilités
(l'\cref{exo:m2:bond-futures:7}).
\end{tutorial}

\section{Construction~: l'analyseur de base}\label{bld:m2:bond-futures:ctd}

\begin{build}
\bfield{Rôle} La firme miniature couvre ses positions obligataires avec des contrats à terme
et mène des \omterm{def:m2:bond-futures:irr}{basis trades}~; les deux ont besoin du facteur, du \omterm{def:m2:bond-futures:cf}{prix de facturation}, de la base et
du \omterm{def:m2:bond-futures:irr}{taux de repo implicite} de chaque titre livrable, chaque jour. \bfield{Interface} \texttt{conversion\_factor(coupon, first\_day, maturity,
quarter\_rounding)}~; \texttt{invoice\_price(F, cf, accrued)}~; \texttt{Deliverable(name, clean, accrued\_now, accrued\_delivery, cf,
coupon\_income)}~; \texttt{gross\_basis}~; \texttt{carry\_to\_delivery}~;
\texttt{net\_basis}~; \texttt{implied\_repo}~; \texttt{cheapest\_to\_deliver}.
\bfield{Règles} Facteurs selon la formule de la bourse~; prix pour 100~; financement en
actual/360~; coupon couru issu de la bibliothèque obligataire
(la \cref{bld:m2:government-bonds:bond}).
\bfield{Tests d'acceptation} \texttt{code/firm/ctd/tests/}~: les cinq
facteurs publiés~; comptage en mois entiers~; une note à 6~\% arrivant à échéance le premier jour du
mois de livraison a un facteur de 1~; la \omterm{def:m2:bond-futures:irr}{base nette} du CTD est nulle au
\omterm{def:m2:bond-futures:irr}{taux de repo implicite} du CTD, et celle de toutes les autres notes est positive.
\bfield{Pour aller plus loin} La formule de facteur d'Eurex (coupons annuels)~; le choix de la
date de livraison, issu de l'option de timing~; la valeur de l'\omterm{def:m2:bond-futures:options}{option de qualité} à partir
d'un modèle de courbe plutôt que de déplacements parallèles.
\end{build}

\begin{omsources}
\item CBOT, soumission 25-099 à la CFTC (contract grade of the ten-year note
future), février 2025~; CME Group, \emph{Calculating U.S. Treasury Futures
Conversion Factors}.
\item M.\ Henrard, \emph{Bond Futures: Description and Pricing}, OpenGamma
Quantitative Research, 2011.
\item Federal Reserve, FEDS Notes, «~Decomposing Hedge Funds' U.S. Treasury
Exposures~», juin 2026~; Office of Financial Research, «~Hedge Funds' Cash
Treasury Holdings Reach \$2 Trillion~», août 2026.
\item G.\ Burghardt et T.\ Belton, \emph{The Treasury Bond Basis}, 3e éd.,
McGraw-Hill, 2005.
\end{omsources}

\section{Exercices}

\begin{exercise}[$\star$]\label{exo:m2:bond-futures:1}
Le contrat à terme se règle à 111-24+. Le \omterm{def:m2:bond-futures:cf}{facteur de concordance} du CTD vaut 0.8870 et son
coupon couru à la livraison vaut 1.453125. Donner le \omterm{def:m2:bond-futures:cf}{prix de facturation} pour 100 et par
contrat.
\end{exercise}

\begin{exercise}[$\star$]\label{exo:m2:bond-futures:2}
Le prix pied de coupon du CTD vaut 99.132. Donner sa \omterm{def:m2:bond-futures:irr}{base brute} en prix et en 32es.
Pourquoi une \omterm{def:m2:bond-futures:irr}{base brute} peut-elle être négative~?
\end{exercise}

\begin{exercise}[$\star$]\label{exo:m2:bond-futures:3}
Calculer le facteur du contrat de décembre 2026 d'un billet à 4~\% échéant le 15 novembre 2033, et
d'un billet à 6~\% de même échéance. Pourquoi le second n'est-il pas exactement égal à 1~?
\end{exercise}

\begin{exercise}[$\star\star$]\label{exo:m2:bond-futures:4}
Le \omterm{def:m2:bond-futures:irr}{taux de repo implicite} du CTD vaut 3.82~\% et le \omterm{def:m2:repo-and-specials:rate}{taux de repo} 3.90~\%. On achète la base
pour un nominal de 100~millions de dollars et on la garde jusqu'à la livraison sans basculement. Que perd-on~?
Pourquoi pourrait-on tout de même le faire~?
\end{exercise}

\begin{exercise}[$\star\star$]\label{exo:m2:bond-futures:5}
Un fonds détient pour 100~millions de dollars de nominal du CTD et veut le couvrir avec le
contrat à terme. Combien de contrats doit-il vendre~? Pourquoi la réponse fait-elle intervenir le
\omterm{def:m2:bond-futures:cf}{facteur de concordance}~?
\end{exercise}

\begin{exercise}[$\star\star$]\label{exo:m2:bond-futures:6}
Expliquer à partir de la duration pourquoi le billet le plus court est le \omterm{def:m2:bond-futures:ctd}{moins cher à livrer} lorsque
les rendements sont inférieurs à 6~\%, et pourquoi le basculement se fait vers le billet le plus long quand les
rendements montent.
\end{exercise}

\begin{exercise}[$\star\star\star$]\label{exo:m2:bond-futures:7}
\emph{Programmation.} Avec \texttt{futures\_demo.switch\_option\_value}, valoriser
l'option de basculement pour une volatilité normale de 80 points de base par an sur les 97
jours, puis trois fois cette volatilité. Donner les deux valeurs en 64es et commenter.
\end{exercise}

\begin{exercise}[$\star\star\star$]\label{exo:m2:bond-futures:8}
\emph{Trouver l'erreur.} «~Le contrat à terme à dix ans se traite à 111 alors que les billets qu'il
livre se traitent autour de 100~: le contrat à terme est cher de 11~\%, il suffit de le vendre et d'acheter les
billets.» Corriger le raisonnement.
\end{exercise}

\section{Problème~: Le basculement}

\begin{problem}[{Problème du week-end --- quand le moins cher à livrer change}]\label{pb:m2:bond-futures:1}
Reprendre le panier et les prix de l'\cref{ex:m2:bond-futures:basket}~: le contrat à terme
à 111-24+, \omterm{def:m2:repo-and-specials:rate}{taux de repo} 3.90~\%, 97 jours avant la livraison.

\medskip\noindent\textbf{Partie I --- Le panier.}
\begin{enumerate}
\item Quels billets sont livrables dans le contrat de décembre, selon la règle de
l'\cref{dat:m2:bond-futures:contracts}~? Vérifier chacun des cinq.
\item Donner le facteur du billet à 4~\% d'août 2034 et celui du CTD.
\item Donner la \omterm{def:m2:bond-futures:irr}{base brute}, le carry et la \omterm{def:m2:bond-futures:irr}{base nette} du CTD.
\item Donner son \omterm{def:m2:bond-futures:irr}{taux de repo implicite}. De combien est-il inférieur au \omterm{def:m2:repo-and-specials:rate}{taux de repo}~?
\item Pourquoi son carry est-il négatif~?
\end{enumerate}

\medskip\noindent\textbf{Partie II --- Le basculement.}
\begin{enumerate}[resume]
\item À la livraison, les rendements restant inchangés, quel est le prix de chaque billet par
unité de facteur, par rapport à celui du CTD~?
\item De combien tous les rendements doivent-ils monter pour que le CTD change~? Quel billet prend
le relais~?
\item Quel rendement a le CTD à ce moment-là~?
\item Pourquoi le basculement se produit-il sous 6~\%~?
\item Après le basculement, comment la duration du contrat à terme évolue-t-elle~?
\end{enumerate}

\medskip\noindent\textbf{Partie III --- L'option.}
\begin{enumerate}[resume]
\item Avec une volatilité normale de 80 points de base par an, quel est l'écart-type
des rendements sur les 97 jours~?
\item À combien d'écarts types se trouve le basculement~?
\item Donner la valeur de l'option de basculement à cette volatilité, en 64es.
\item Et à trois fois cette volatilité~?
\item Dans l'exemple, le contrat à terme se traite un 64e sous le forward du CTD. À quoi
d'autre ce 64e pourrait-il rémunérer~?
\end{enumerate}

\medskip\noindent\textbf{Partie IV --- Discernement.}
\begin{enumerate}[resume]
\item Une trader est longue de la base du CTD. Pourquoi est-elle longue en volatilité~?
\item Pourquoi le basculement compte-t-il pour un hedger qui utilise des contrats à terme, et non seulement pour un
basis trader~?
\item Que ferait la réouverture du billet d'août 2033, qui accroît son offre, sur
la base~?
\item Énoncer le \emph{résultat nommé}~: le déplacement à partir duquel le CTD bascule, et
la valeur de l'option de basculement pour les deux volatilités.
\item En une phrase~: que livre réellement un contrat à terme sur obligations~?
\end{enumerate}
\end{problem}

\section{Questions d'entretien}

\begin{interviewq}[$\star$ \iqroles{trader, banque}]\label{iq:m2:bond-futures:1}
Qu'est-ce qu'un \omterm{def:m2:bond-futures:cf}{facteur de concordance}, et pourquoi un contrat à terme sur Treasuries en a-t-il besoin~?
\end{interviewq}

\begin{interviewq}[$\star$ \iqroles{trader, chercheur}]\label{iq:m2:bond-futures:2}
Comment identifiez-vous le \omterm{def:m2:bond-futures:ctd}{moins cher à livrer}~?
\end{interviewq}

\begin{interviewq}[$\star\star$ \iqroles{trader, chercheur}]\label{iq:m2:bond-futures:3}
Qu'est-ce que le \omterm{def:m2:bond-futures:irr}{taux de repo implicite}, et que signifie-t-il si le \omterm{def:m2:bond-futures:irr}{taux de repo implicite} du CTD
est supérieur au \omterm{def:m2:repo-and-specials:rate}{taux de repo} effectif~?
\end{interviewq}

\begin{interviewq}[$\star\star$ \iqroles{trader, banque}]\label{iq:m2:bond-futures:4}
Expliquez le \omterm{def:m2:bond-futures:irr}{basis trade} sur Treasuries et ce qui peut mal tourner.
\end{interviewq}

\begin{interviewq}[$\star\star$ \iqroles{chercheur, développeur}]\label{iq:m2:bond-futures:5}
Comment calculeriez-vous le DV01 d'un contrat à terme sur Treasuries~?
\end{interviewq}

\begin{interviewq}[$\star\star\star$ \iqroles{chercheur, trader}]\label{iq:m2:bond-futures:6}
Quelles options le vendeur d'un contrat à terme sur Treasuries détient-il, et comment
valoriseriez-vous la plus importante~?
\end{interviewq}
